\setcounter{section}{4}

  \zwischenueberschrift{Übungsaufgaben}

\inputaufgabe
{}
{

Es sei

\mavergleichskette
{\vergleichskette
{ h(x_1 , \ldots , x_n)
}
{ = }{ \sum_{i = 1}^n a_ix_i
}
{  }{
}
{    }{
}
{   }{
}
}
{}{}{}
eine
\definitionsverweis {Linearform}{}{}
$\neq 0$ und

\mavergleichskette
{\vergleichskette
{ Y
}
{ = }{ h^{-1}(c)
}
{  }{
}
{    }{
}
{   }{
}
}
{}{}{}
die zugehörige Hyperebene. Zeige, dass die
\definitionsverweis {Weingartenabbildung}{}{}
die Nullabbildung ist.

}
{} {}

\inputaufgabe
{}
{

Es sei

\mavergleichskettedisp
{\vergleichskette
{ Y
}
{ =} { \left\{  \left(  x   , \,   y  , \,  z                 \right)   \mid   x^2+y^2 = z^2 , \,  \left(  x   , \,   y  , \,  z                 \right)  \neq  \left(  0   , \,   0  , \,   0                 \right)       \right\}
}
{ \subseteq} { W
}
{ =} { \R^3 \setminus \{  \left(  0   , \,   0  , \,   0                 \right) \}
}
{ } {
}
}
{}{}{}
der sogenannte Standardkegel und sei

\mavergleichskette
{\vergleichskette
{ P
}
{ = }{ \left(  x   , \,   y  , \,  z                 \right)
}
{ \in }{ Y
}
{    }{
}
{   }{
}
}
{}{}{.}
Es sei

\mavergleichskettedisp
{\vergleichskette
{ \gamma(t)
}
{ =} { \left(  x  \cos  t    , \,   y  \sin  t  , \,  z                 \right)
}
{ } {
}
{ } {
}
{ } {
}
}
{}{}{.}
Bestimme

\mathdisp {\operatorname{lim}_{   t  \rightarrow  0   } \,  { \frac{  N( \gamma(t) ) -  N( \gamma(0) )  }{ t } }} { , }

wobei $N$ das Einheitsnormalenvektorfeld zum Gradienten von $x^2+y^2-z^2$ sei.

}
{} {}

\inputaufgabe
{}
{

Es sei

\mavergleichskettedisp
{\vergleichskette
{ Y
}
{ =} { \left\{  \left(  x   , \,   y  , \,  z                 \right)   \mid   x^2+y^2 =1         \right\}
}
{ \subseteq} { \R^3
}
{ } {
}
{ } {
}
}
{}{}{}
der sogenannte Standardzylinder und sei

\mavergleichskette
{\vergleichskette
{ P
}
{ = }{ \left(  1   , \,   0  , \,  z                 \right)
}
{ \in }{ Y
}
{    }{
}
{   }{
}
}
{}{}{.}
Es sei

\mavergleichskettedisp
{\vergleichskette
{ \gamma(t)
}
{ =} { \left(  \cos  t   , \,   \sin  t  , \,  z                 \right)
}
{ } {
}
{ } {
}
{ } {
}
}
{}{}{.}
Bestimme

\mathdisp {\operatorname{lim}_{   t  \rightarrow  0   } \,  { \frac{  N( \gamma(t) ) -  N( \gamma(0) )  }{ t } }} { , }

wobei $N$ das Einheitsnormalenvektorfeld zum Gradienten von $x^2+y^2-1$ sei.

}
{} {}

\inputaufgabe
{}
{

Es sei

\mavergleichskettedisp
{\vergleichskette
{ Y
}
{ =} { \left\{  \left(  x   , \,   y  , \,  z                 \right)   \mid   x^2+y^2 =1         \right\}
}
{ \subseteq} { \R^3
}
{ } {
}
{ } {
}
}
{}{}{}
und sei

\mavergleichskette
{\vergleichskette
{ P
}
{ \in }{ Y
}
{  }{
}
{    }{
}
{   }{}
}
{}{}{}
ein Punkt. Zeige, dass die
\definitionsverweis {Weingartenabbildung}{}{}
$L_P$ nicht bijektiv, aber auch nicht die Nullabbildung ist.

}
{} {}

\inputaufgabe
{}
{

Zeige, dass eine
\definitionsverweis {differenzierbare Fläche}{}{}

\mavergleichskette
{\vergleichskette
{ Y
}
{ \subseteq }{ \R^3
}
{  }{
}
{    }{
}
{   }{
}
}
{}{}{}
eine Gerade

\mavergleichskette
{\vergleichskette
{ G
}
{ \subseteq }{ Y
}
{  }{
}
{    }{
}
{   }{
}
}
{}{}{}
enthalten kann, die, aufgefasst im
\definitionsverweis {Tangentialraum}{}{}
$T_PY$ für einen Punkt

\mavergleichskette
{\vergleichskette
{ P
}
{ \in }{ G
}
{  }{
}
{    }{
}
{   }{
}
}
{}{}{,}
nicht zum
\definitionsverweis {Kern}{}{}
der
\definitionsverweis {Weingartenabbildung}{}{}
$L_P$ gehören muss.

}
{} {}

\inputaufgabe
{}
{

Es sei

\mavergleichskette
{\vergleichskette
{ W
}
{ \subseteq }{ \R^n
}
{  }{
}
{    }{
}
{   }{
}
}
{}{}{}
\definitionsverweis {offen}{}{,}
\maabb {h} { W } { \R
} {}
eine
\definitionsverweis {stetig differenzierbare Funktion}{}{}
und

\mavergleichskette
{\vergleichskette
{ Y
}
{ = }{ h^{-1}(c)
}
{  }{
}
{    }{
}
{   }{
}
}
{}{}{}
die
\definitionsverweis {Faser}{}{}
zu

\mavergleichskette
{\vergleichskette
{ c
}
{ \in }{ \R
}
{  }{
}
{    }{
}
{   }{
}
}
{}{}{,}
wobei $h$ in jedem Punkt von $Y$
\definitionsverweis {regulär}{}{}
sei. Wir fassen $h$ auch als eine Abbildung $\tilde{h}$  auf
\mathl{W \times \R^m}{} auf, wobei die hinteren $m$ Variablen nicht explizit in $\tilde{h}$ eingehen. Es sei

\mavergleichskettedisp
{\vergleichskette
{ Z
}
{ =} { \tilde{h}^{-1}(c)
}
{ =} { Y \times \R^m
}
{ \subseteq} { W \times \R^{m}
}
{ } {
}
}
{}{}{.}
Es sei

\mavergleichskette
{\vergleichskette
{ Q
}
{ \in }{ Z
}
{  }{
}
{    }{
}
{   }{
}
}
{}{}{}
ein Punkt über

\mavergleichskette
{\vergleichskette
{ P
}
{ \in }{ Y
}
{  }{
}
{    }{
}
{   }{
}
}
{}{}{.}
In welcher Beziehung stehen die
\definitionsverweis {Weingartenabbildungen}{}{}
\mathkor {} {L_Q} {und} {L_P} {?}

}
{} {}

\inputaufgabegibtloesung
{}
{

Es seien
\mathkor {} {r} {und} {R} {}
positive reelle Zahlen mit

\mavergleichskette
{\vergleichskette
{ 0
}
{ < }{ r
}
{ < }{ R
}
{    }{
}
{   }{
}
}
{}{}{,}
wir betrachten den
\zusatzklammer {eingebetteten} {} {}
Torus

\mavergleichskettedisp
{\vergleichskette
{ T
}
{ =} { \left\{ (x,y,z) \in \R^3  \mid   \left(  \sqrt{x^2+y^2} -R  \right)^2 +z^2  = r^2         \right\}
}
{ } {
}
{ } {
}
{ } {
}
}
{}{}{.}
\aufzaehlungdrei{Bestimme ein
\definitionsverweis {Einheitsnormalenfeld}{}{}
$N$ zu $T$.
}{Bestimme im Punkt
\mathl{\left( R-r  , \,   0  , \,   0                 \right)}{} und für den Vektor

\mavergleichskette
{\vergleichskette
{ v
}
{ = }{ \begin{pmatrix}  0  \\ 1 \\  0   \end{pmatrix}
}
{  }{
}
{    }{
}
{   }{
}
}
{}{}{}
die
\definitionsverweis {Richtungsableitung}{}{}

\mathl{D_vN}{.}
}{Bestimme im Punkt
\mathl{\left( R-r  , \,   0  , \,   0                 \right)}{} und für den Vektor

\mavergleichskette
{\vergleichskette
{ w
}
{ = }{ \begin{pmatrix}  0  \\ 0\\  1   \end{pmatrix}
}
{  }{
}
{    }{
}
{   }{
}
}
{}{}{}
die
\definitionsverweis {Richtungsableitung}{}{}

\mathl{D_w N}{.}
}

}
{} {}

\inputaufgabe
{}
{

Bestimme für den Graphen $Y$  der Funktion

\mavergleichskettedisp
{\vergleichskette
{ f(x,y)
}
{ =} { xy
}
{ } {
}
{ } {
}
{ } {
}
}
{}{}{}
die
\definitionsverweis {Weingartenabbildung}{}{}
$L_P$ für jeden Punkt

\mavergleichskettedisp
{\vergleichskette
{ P
}
{ =} { (x,y,f(x,y) )
}
{ } {
}
{ } {
}
{ } {
}
}
{}{}{.}
Bestimme die Eigenwerte, die Eigenvektoren und eine Diagonalmatrix für $L_P$.

}
{} {}

\inputaufgabe
{}
{

Es sei

\mavergleichskette
{\vergleichskette
{ W
}
{ \subseteq }{ \R^n
}
{  }{
}
{    }{
}
{   }{
}
}
{}{}{}
\definitionsverweis {offen}{}{,}
\maabb {h} { W } { \R
} {}
eine
\definitionsverweis {stetig differenzierbare Funktion}{}{}
und

\mavergleichskette
{\vergleichskette
{ Y
}
{ = }{ h^{-1}(c)
}
{  }{
}
{    }{
}
{   }{
}
}
{}{}{}
die
\definitionsverweis {Faser}{}{}
zu

\mavergleichskette
{\vergleichskette
{ c
}
{ \in }{ \R
}
{  }{
}
{    }{
}
{   }{
}
}
{}{}{,}
wobei $h$ in jedem Punkt von $Y$
\definitionsverweis {regulär}{}{}
sei. Es sei
\maabbdisp {L} {\R^n } { \R^n
} {}
eine
\definitionsverweis {lineare Isometrie}{}{,}

\mavergleichskette
{\vergleichskette
{ W'
}
{ = }{ L^{-1}(W)
}
{  }{
}
{    }{
}
{   }{
}
}
{}{}{}
und

\mavergleichskettedisp
{\vergleichskette
{ Z
}
{ =} { L^{-1}(Y)
}
{ =} { (h \circ L)^{-1}(c)
}
{ } {
}
{ } {
}
}
{}{}{.}
Zeige, dass sich die
\definitionsverweis {Weingartenabbildung}{}{}
$L_P$ zu

\mavergleichskette
{\vergleichskette
{ P
}
{ \in }{ Y
}
{  }{
}
{    }{
}
{   }{
}
}
{}{}{}
und die Weingartenabbildung zu

\mavergleichskette
{\vergleichskette
{ Q
}
{ = }{ L^{-1}(P)
}
{ \in }{ Z
}
{    }{
}
{   }{
}
}
{}{}{}
entsprechen, wenn man mithilfe von $L$ die Tangentialräume $T_PY$ und $T_QZ$ miteinander in Beziehung setzt.

}
{} {}

\inputaufgabegibtloesung
{}
{

Bestimme für die durch

\mavergleichskettedisp
{\vergleichskette
{ x^2+y^2+2z^2
}
{ =} { 4
}
{ } {
}
{ } {
}
{ } {
}
}
{}{}{}
gegebene Fläche $Y$ und den Punkt

\mavergleichskettedisp
{\vergleichskette
{ P
}
{ =} { \left(  1  , \,   1  , \,   1                 \right)
}
{ \in} { Y
}
{ } {
}
{ } {
}
}
{}{}{}
eine Diagonalmatrix für die
\definitionsverweis {Weingartenabbildung}{}{}
$L_P$.

}
{} {}

\inputaufgabe
{}
{

Es sei

\mavergleichskette
{\vergleichskette
{ W
}
{ \subseteq }{ \R^n
}
{  }{
}
{    }{
}
{   }{
}
}
{}{}{}
\definitionsverweis {offen}{}{,}
\maabb {h} { W } { \R
} {}
eine zweifach
\definitionsverweis {stetig differenzierbare Funktion}{}{}
und

\mavergleichskette
{\vergleichskette
{ Y
}
{ = }{ h^{-1}(c)
}
{  }{
}
{    }{
}
{   }{
}
}
{}{}{}
die
\definitionsverweis {Faser}{}{}
zu

\mavergleichskette
{\vergleichskette
{ c
}
{ \in }{ \R
}
{  }{
}
{    }{
}
{   }{
}
}
{}{}{,}
wobei $h$ in jedem Punkt von $Y$
\definitionsverweis {regulär}{}{}
sei. Wie ändert sich die
\definitionsverweis {Weingartenabbildung}{}{}
und wie die zugehörigen
\definitionsverweis {Eigenräume}{}{,}
wenn man von einer
\definitionsverweis {Orientierung}{}{}
von $Y$ zur entgegengesetzten Orientierung wechselt?

}
{} {}

  \zwischenueberschrift{Aufgaben zum Abgeben}

\inputaufgabe
{4}
{

Bestimme in
[[Einschaliges Hyperboloid/Weingartenabbildung/Beispiel|Beispiel 4.5]]
die beschreibende Matrix für die
\definitionsverweis {Weingartenabbildung}{}{}
für Punkte der Form
\mathl{\begin{pmatrix}  0  \\ y \\ z   \end{pmatrix}}{} bezüglich der Basis
\mathkor {} {\begin{pmatrix}  y  \\ 0 \\  0   \end{pmatrix}} {und} {\begin{pmatrix}  0  \\ z \\  y   \end{pmatrix}} {}
des Tangentialraumes. Bestimme die Eigenwerte und eine Basis aus Eigenvektoren.

}
{} {}

\inputaufgabe
{5}
{

Bestimme für die durch

\mavergleichskettedisp
{\vergleichskette
{ x^2+2y^2+5z^2
}
{ =} { 11
}
{ } {
}
{ } {
}
{ } {
}
}
{}{}{}
gegebene Fläche $Y$ und den Punkt

\mavergleichskettedisp
{\vergleichskette
{ P
}
{ =} { \left(  2  , \,   1  , \,   1                 \right)
}
{ \in} { Y
}
{ } {
}
{ } {
}
}
{}{}{}
eine Diagonalmatrix für die
\definitionsverweis {Weingartenabbildung}{}{}
$L_P$.

}
{} {}

\inputaufgabe
{4}
{

Bestimme für die durch

\mavergleichskettedisp
{\vergleichskette
{ x^2+2y^2+3z^2+4w^2
}
{ =} { 10
}
{ } {
}
{ } {
}
{ } {
}
}
{}{}{}
gegebene Hyperfläche

\mavergleichskette
{\vergleichskette
{ Y
}
{ \subseteq }{ \R^4
}
{  }{
}
{    }{
}
{   }{
}
}
{}{}{}
und den Punkt

\mavergleichskettedisp
{\vergleichskette
{ P
}
{ =} { \left(  1  , \,   1  , \,   1 , \,   1                \right)
}
{ \in} { Y
}
{ } {
}
{ } {
}
}
{}{}{}
eine Matrix für die
\definitionsverweis {Weingartenabbildung}{}{}
$L_P$ bezüglich einer geeigneten
\definitionsverweis {Basis}{}{}
des
\definitionsverweis {Tangentialraumes}{}{}
$T_PY$.

}
{} {}

\inputaufgabe
{6}
{

Bestimme für den Graphen $Y$ der Funktion

\mavergleichskettedisp
{\vergleichskette
{ f(x,y)
}
{ =} { x^3y^5
}
{ } {
}
{ } {
}
{ } {
}
}
{}{}{}
die
\definitionsverweis {Weingartenabbildung}{}{}
$L_P$ für die Punkte

\mavergleichskette
{\vergleichskette
{ P_1
}
{ = }{ (0,0,0)
}
{  }{
}
{    }{
}
{   }{
}
}
{}{}{,}

\mavergleichskette
{\vergleichskette
{ P_2
}
{ = }{ (0,1,0)
}
{  }{
}
{    }{
}
{   }{
}
}
{}{}{,}

\mavergleichskette
{\vergleichskette
{ P_3
}
{ = }{ (2,1,8)
}
{  }{
}
{    }{
}
{   }{
}
}
{}{}{.}
Bestimme jeweils die Eigenwerte, die Eigenvektoren und eine Diagonalmatrix für $L_P$.

}
{} {}

 [[Kategorie:Latexseite]]
